\section{Dissipation generated by two Pauli matrices}
\label{sec:two_pauli_dissipation}
%% ---------------------------------------------------------

\begin{theorem}[Pauli conjugation identities]
    \label{thm:pauli_conjugation_identities}
    \lean{PauliDissipation.pauli_conjTranspose}
    \lean{PauliDissipation.pauli_sq}
    \lean{PauliDissipation.pauli_mul_pauli_mul}
    \lean{PauliDissipation.dissipator_pauli}
    \leanok
    For Pauli directions $a,k\in\{x,y,z\}$,
    \begin{align}
        \sigma_a^\dagger&=\sigma_a,
        \notag\\
        \sigma_a^2&=\Id,
        \notag\\
        \sigma_a\sigma_k\sigma_a
        &=
        \begin{cases}
            \sigma_k,&a=k,\\
            -\sigma_k,&a\ne k.
        \end{cases}
        \notag
    \end{align}
    Consequently, the dissipator generated by $\sigma_a$ is
    $\mathcal D_{\sigma_a}(X)=\sigma_aX\sigma_a-X$.
\end{theorem}

\begin{proof}
    \leanok
    Direct multiplication of the three Pauli matrices proves the identities;
    substituting $\sigma_a^\dagger=\sigma_a$ and $\sigma_a^2=\Id$ into the
    Lindblad dissipator gives the final formula.
\end{proof}

\begin{definition}[Two-Pauli dissipative generator]
    \label{def:two_pauli_generator}
    \lean{PauliDissipation.twoPauliGenerator}
    \lean{PauliDissipation.twoPauliGenerator_apply}
    \lean{PauliDissipation.twoPauliGenerator_one}
    \leanok
    For Pauli directions $a,b\in\{x,y,z\}$ and real rates
    $\gamma_a,\gamma_b$, define
    \begin{align}
        \mathcal L_{a,b}(X)
        &=
        \gamma_a(\sigma_aX\sigma_a-X)
        +\gamma_b(\sigma_bX\sigma_b-X).
        \label{eq:semigroup_two_pauli_generator}
    \end{align}
    In particular, $\mathcal L_{a,b}(\Id)=0$.
\end{definition}

\begin{theorem}[Pauli spectrum of the two-direction generator]
    \label{thm:two_pauli_generator_spectrum}
    \lean{PauliDissipation.twoPauliGenerator_pauli}
    \lean{PauliDissipation.twoPauliGenerator_pauli_a}
    \lean{PauliDissipation.twoPauliGenerator_pauli_b}
    \lean{PauliDissipation.twoPauliGenerator_pauli_remaining}
    \leanok
    \uses{def:two_pauli_generator}
    For Pauli directions $a,b,k\in\{x,y,z\}$ and real rates
    $\gamma_a,\gamma_b$,
    \begin{align}
        \mathcal L_{a,b}(\sigma_k)
        &=-2\bigl(\gamma_a\mathbf 1_{a\ne k}
                  +\gamma_b\mathbf 1_{b\ne k}\bigr)\sigma_k.
        \label{eq:semigroup_two_pauli_spectrum}
    \end{align}
    Hence, when $a\ne b$, the eigenvalues on $\sigma_a$, $\sigma_b$, and the
    remaining Pauli matrix are respectively
    $-2\gamma_b$, $-2\gamma_a$, and $-2(\gamma_a+\gamma_b)$.
\end{theorem}

\begin{proof}
    \leanok
    \uses{thm:pauli_conjugation_identities}
    Apply Theorem~\ref{thm:pauli_conjugation_identities} separately to the two
    dissipators in (\ref{eq:semigroup_two_pauli_generator}) and collect the two
    scalar coefficients.
\end{proof}

\begin{theorem}[Trace symmetry of two-Pauli dissipation]
    \label{thm:two_pauli_generator_trace_symmetry}
    \lean{PauliDissipation.dissipator_pauli_trace_symmetry}
    \lean{PauliDissipation.twoPauliGenerator_trace_conjTranspose_symmetry}
    \leanok
    \uses{def:two_pauli_generator}
    For Pauli directions $a,b\in\{x,y,z\}$, real rates
    $\gamma_a,\gamma_b$, and all $X,Y\in\MN{2}$,
    \begin{align}
        \tr(X^\dagger\mathcal L_{a,b}(Y))
        &=\tr(\mathcal L_{a,b}(X)^\dagger Y).
        \notag
    \end{align}
\end{theorem}

\begin{proof}
    \leanok
    \uses{thm:pauli_conjugation_identities}
    Cyclicity of the trace makes each map
    $X\mapsto\sigma_aX\sigma_a-X$ symmetric for the Hilbert--Schmidt pairing.
    Real linear combinations preserve this symmetry.
\end{proof}

\begin{theorem}[Scalar kernel of two-Pauli dissipation]
    \label{thm:two_pauli_generator_scalar_kernel}
    \lean{PauliDissipation.twoPauliGenerator_eq_zero_iff}
    \lean{PauliDissipation.pauli_sum_eq_zero_iff}
    \leanok
    \uses{def:two_pauli_generator}
    Suppose $a\ne b$ and $\gamma_a,\gamma_b>0$. For $X\in\MN{2}$,
    $\mathcal L_{a,b}(X)=0$ if and only if $X\in\C\Id$.
\end{theorem}

\begin{proof}
    \leanok
    \uses{thm:two_pauli_generator_spectrum}
    Expand $X$ in the basis $\{\Id,\sigma_x,\sigma_y,\sigma_z\}$. The three
    Pauli coefficients vanish because the three eigenvalues in
    (\ref{eq:semigroup_two_pauli_spectrum}) are nonzero, while the identity
    coefficient is unrestricted.
\end{proof}

\begin{definition}[Maximally mixed qubit state]
    \label{def:two_pauli_maximally_mixed}
    \lean{PauliDissipation.maximallyMixed}
    \leanok
    The maximally mixed qubit state is $\tau_1=\frac12\Id_2$.
\end{definition}

\begin{theorem}[Faithful stationary state of two-Pauli dissipation]
    \label{thm:two_pauli_maximally_mixed_fixed}
    \lean{PauliDissipation.twoPauliGenerator_hasSimpleFaithfulKernel}
    \lean{PauliDissipation.maximallyMixed_posDef}
    \lean{PauliDissipation.maximallyMixed_mem_densityMatrices}
    \lean{PauliDissipation.twoPauliGenerator_maximallyMixed}
    \lean{PauliDissipation.twoPauliGenerator_hasSimpleKernel}
    \lean{PauliDissipation.expSemigroup_fixed_iff_scalar}
    \lean{PauliDissipation.autonomousSemigroup_fixed_iff_scalar}
    \lean{PauliDissipation.expSemigroup_fixes_maximallyMixed}
    \leanok
    \uses{def:two_pauli_generator, def:two_pauli_maximally_mixed,
        def:exp_semigroup_ch12, def:density_matrices}
    Suppose $a\ne b$ and $\gamma_a,\gamma_b>0$. The state $\tau_1$ is positive
    definite and belongs to $\mc{D}_2$. Furthermore,
    $\mathcal L_{a,b}(\tau_1)=0$, the kernel of $\mathcal L_{a,b}$ is the
    one-dimensional space $\C\tau_1$, and
    $e^{t\mathcal L_{a,b}}(\tau_1)=\tau_1$ for every $t\ge0$. More generally,
    for $X\in\MN{2}$,
    \begin{align}
        (\forall t\ge0,\ e^{t\mathcal L_{a,b}}(X)=X)
        &\quad\Longleftrightarrow\quad
        X\in\C\Id.
        \notag
    \end{align}
    The same equivalence holds for any autonomous semigroup $T_t$ satisfying
    $T_t=e^{t\mathcal L_{a,b}}$ at every non-negative time.
\end{theorem}

\begin{proof}
    \leanok
    \uses{thm:two_pauli_generator_scalar_kernel}
    The defining formula gives $\mathcal L_{a,b}(\tau_1)=0$.
    Theorem~\ref{thm:two_pauli_generator_scalar_kernel} identifies its kernel,
    and exponentiating a kernel vector leaves it fixed. Conversely, a vector
    fixed for every non-negative time lies in the generator kernel.
\end{proof}

\begin{theorem}[Conditional logarithmic-time trace-norm estimate]
    \label{thm:two_pauli_conditional_trace_norm_mixing}
    \lean{PauliDissipation.q3_traceNorm_le_rpow_of_logarithmic_time}
    \lean{PauliDissipation.squared_trace_distance_bound_of_entropy_decay_and_pinsker}
    \lean{PauliDissipation.squared_trace_distance_le_rpow_of_exp_threshold}
    \lean{PauliDissipation.q3_traceNorm_sq_bound_of_entropy_decay_and_pinsker}
    \lean{PauliDissipation.q3_exp_threshold_of_logarithmic_time}
    \leanok
    \uses{def:trace_norm, def:quantum_relative_entropy}
    Let $N\in\N$ with $N\ge1$, let $p\ge0$ and $\gamma_*>0$, and let
    $\rho_t=\Phi_t(\rho)$ for matrices $\rho,\omega\in\MN{2^N}$.
    Assume at the time under consideration that
    \begin{align}
        D(\rho_t\Vert\omega)
            &\le e^{-(\gamma_*/2)t}D(\rho\Vert\omega),
            \notag\\
        D(\rho\Vert\omega)&\le N\log 2,
            \notag\\
        \lVert\rho_t-\omega\rVert_1^2
            &\le 2D(\rho_t\Vert\omega).
            \notag
    \end{align}
    If
    \begin{align}
        t&\ge \frac{2}{\gamma_*}
        ((2p+1)\log N+\log(2\log2)),
        \label{eq:semigroup_mixing_time}
    \end{align}
    then $\lVert\rho_t-\omega\rVert_1\le N^{-p}$.
    The complete modified logarithmic Sobolev estimate motivating the first
    hypothesis is \cite[Theorems~1.1 and~3.3]{GaoRouze2022Complete}; that
    estimate, its tensorization, and quantum Pinsker remain hypotheses here.
\end{theorem}

\begin{proof}
    \leanok
    Entropy decay and Pinsker first give
    \begin{align}
        \lVert\rho_t-\omega\rVert_1^2
        &\le 2N\log2\,e^{-\gamma_*t/2}.
        \notag
    \end{align}
    The time hypothesis (\ref{eq:semigroup_mixing_time}) implies
    $e^{-\gamma_*t/2}\le N^{-2p}/(2N\log2)$, so the squared trace norm is at
    most $N^{-2p}$. Both sides are non-negative, and taking square roots gives
    the result.
\end{proof}

%% ---------------------------------------------------------
\section{Primitivity and irreducibility of QDS}
\label{sec:qds_primitivity}
%% ---------------------------------------------------------

%% Source: \cite[Section~7.1, Proposition~7.5]{Wolf2012Quantum}.
\subsection{Auxiliary spectral and semigroup lemmas}

\begin{definition}[Quantum dynamical semigroup]
    \label{def:is_quantum_dyn_semigroup_ch12}
    \lean{IsQuantumDynSemigroup}\leanok
    \uses{def:continuous_dyn_semigroup_ch12, def:channel}
    A quantum dynamical semigroup is a norm-continuous dynamical semigroup whose
    time slices are quantum channels for every $t \ge 0$.
\end{definition}

\begin{theorem}[Semigroup power identity]
    \lean{semigroup_pow}\leanok
    \uses{def:exp_semigroup_ch12}
    For $t\ge 0$ and $n\in\N$, one has $T_{nt}=(T_t)^n$.
\end{theorem}

\begin{lemma}[Density matrices are nonzero]
    \lean{ne_zero_of_mem_densityMatrices'}\leanok
    Every density matrix is nonzero.
\end{lemma}

\begin{theorem}[Compression invariance is stable under powers]
    \lean{compression_preserved_by_pow}\leanok
    If a corner compression is invariant under $E$, it remains invariant under
    every power $E^n$.
\end{theorem}

\begin{theorem}[Generator eigenvectors along the exponential semigroup]
    \lean{expSemigroup_apply_eigenvector}\leanok
    If $L(X)=\mu X$, then $e^{tL}(X)=e^{t\mu}X$ for all $t\in\R$.
\end{theorem}

\begin{theorem}[Spectral mapping for semigroup generators]
    \lean{eigenvalue_exp_of_eigenvalue_generator}\leanok
    If $\mu$ is in the spectrum of $L$, then $e^{t\mu}$ is in the spectrum
    of $e^{tL}$.
\end{theorem}

\begin{theorem}[Root-of-unity rigidity for phases]
    \lean{eq_zero_of_exp_mul_I_isRootOfUnity}\leanok
    If $e^{it\theta}$ is a root of unity for every $t>0$, then $\theta=0$.
\end{theorem}

\begin{theorem}[Peripheral generator eigenvalues are purely imaginary]
    \lean{re_eq_zero_of_peripheral_generator}\leanok
    Peripheral spectral points of the generator have vanishing real part.
\end{theorem}

\begin{theorem}[Peripheral cardinality bound]
    \lean{peripheral_card_le_finrank}\leanok
    The number of peripheral eigenvalues is bounded by $\dim(\MN{D})$.
\end{theorem}

\begin{theorem}[Bounded order under peripheral power-closure]
    \lean{bounded_root_of_peripheral_closed_powers}\leanok
    If all powers of a peripheral eigenvalue remain peripheral, then some
    positive power not exceeding $\dim(\MN{D})$ equals $1$.
\end{theorem}

\begin{theorem}[Peripheral powers stay peripheral under irreducibility]
    \lean{peripheral_powers_closed_of_irreducible_channel_with_fixed}\leanok
    For an irreducible channel with faithful fixed point, every power of a
    peripheral eigenvalue is again peripheral.
\end{theorem}

\begin{theorem}[Spectral-radius criterion from eigenvalue bounds]
    \lean{spectralRadius_lt_one_of_eigenvalues_lt_one}\leanok
    If all eigenvalues satisfy $|\lambda|<1$, then the spectral radius is $<1$.
\end{theorem}

\begin{theorem}[Primitive powers annihilate trace-zero matrices]
    \label{thm:primitive_channel_pow_tendsto_zero_ch12}
    \lean{primitive_channel_pow_tendsto_zero_of_trace_zero}\leanok
    If $E$ is an irreducible primitive channel with faithful fixed density $\sigma$,
    then $E^n(X)\to 0$ for every matrix $X$ with $\tr(X)=0$.
\end{theorem}

\begin{proof}
    \leanok
    Decompose $E^n$ into the fixed-point projection and its complementary part.
    On the complement, every eigenvalue has modulus $<1$, so the complementary powers
    tend to $0$; the trace-zero hypothesis removes the fixed-point component.
\end{proof}

\begin{theorem}[Fixed points of an irreducible channel are trace multiples]
    \label{thm:fixedPoint_eq_trace_smul_irreducible_channel_ch12}
    \lean{fixedPoint_eq_trace_smul_of_irreducible_channel}\leanok
    If $E$ is an irreducible channel with faithful fixed density $\sigma$, then every
    nonzero fixed point $V$ of $E$ satisfies $V=\tr(V)\sigma$.
\end{theorem}

\begin{proof}
    \leanok
    Subtract the trace multiple $\tr(V)\sigma$ and apply the vanishing theorem for
    trace-zero fixed points of an irreducible channel.
\end{proof}

\begin{corollary}[Nonzero fixed points have nonzero trace]
    \label{cor:trace_ne_zero_fixedpoint_irreducible_channel_ch12}
    \lean{trace_ne_zero_of_nonzero_fixedPoint_of_irreducible_channel}\leanok
    Every nonzero fixed point of an irreducible channel has nonzero trace.
\end{corollary}

\begin{proof}
    \leanok
    Otherwise the preceding theorem would force the fixed point to vanish.
\end{proof}

\begin{theorem}[Peripheral eigenvectors become periodic fixed points]
    \label{thm:peripheral_power_fixed_eigenvector_ch12}
    \lean{exists_power_fixed_eigenvector_of_peripheral}\leanok
    For every positive time $t$, a quantum dynamical semigroup irreducible at every
    positive time has the following property: every peripheral eigenvalue $\mu$ of $T_t$
    admits a nonzero eigenvector $V$ and a positive integer $p$ such that
    $T_t(V)=\mu V$ and $T_{pt}(V)=V$.
\end{theorem}

\begin{proof}
    \leanok
    Peripheral powers remain peripheral for an irreducible channel with faithful fixed point.
    The bounded-order lemma then gives $\mu^p=1$, and the semigroup power identity turns the
    eigenvector equation for $T_t$ into a fixed-point equation for $T_{pt}$.
\end{proof}

\begin{theorem}[Peripheral eigenvectors with nonzero trace]
    \label{thm:peripheral_trace_ne_zero_eigenvector_ch12}
    \lean{exists_trace_ne_zero_eigenvector_of_peripheral}\leanok
    For every positive time $t$, under the same hypotheses every peripheral eigenvalue
    of $T_t$ has a nonzero eigenvector with nonzero trace.
\end{theorem}

\begin{proof}
    \leanok
    \uses{thm:peripheral_power_fixed_eigenvector_ch12}
    Apply Theorem~\ref{thm:peripheral_power_fixed_eigenvector_ch12} and then use the previous
    corollary at the irreducible time $pt$.
\end{proof}

\begin{theorem}[Trace-preserving eigenvectors force eigenvalue $1$]
    \label{thm:eigenvalue_eq_one_trace_preserving_eigenvector_ch12}
    \lean{eigenvalue_eq_one_of_trace_preserving_eigenvector}\leanok
    If $E$ is trace-preserving, $E(V)=\mu V$, and $\tr(V)\ne 0$, then $\mu=1$.
\end{theorem}

\begin{proof}
    \leanok
    Take traces in $E(V)=\mu V$ and use trace preservation.
\end{proof}

\begin{theorem}[Peripheral collapse under irreducibility at all times]
    \label{thm:peripheral_eq_one_irreducible_all_ch12}
    \lean{peripheral_eq_one_of_irreducible_all}\leanok
    If a quantum dynamical semigroup is irreducible at every positive time, then every
    peripheral eigenvalue of every positive-time slice equals $1$.
\end{theorem}

\begin{proof}
    \leanok
    \uses{thm:eigenvalue_eq_one_trace_preserving_eigenvector_ch12,
        thm:peripheral_trace_ne_zero_eigenvector_ch12}
    Combine Theorem~\ref{thm:peripheral_trace_ne_zero_eigenvector_ch12} with
    Theorem~\ref{thm:eigenvalue_eq_one_trace_preserving_eigenvector_ch12}.
\end{proof}

\begin{theorem}[Primitivity from irreducibility at all times]
    \label{thm:primitive_of_irreducible_all_ch12}
    \lean{primitive_of_irreducible_all}\leanok
    If a quantum dynamical semigroup is irreducible at every positive time, then every
    positive-time slice is primitive.
\end{theorem}

\begin{proof}
    \leanok
    \uses{thm:peripheral_eq_one_irreducible_all_ch12}
    The unique faithful fixed density gives the one-dimensional peripheral spectrum criterion,
    and Theorem~\ref{thm:peripheral_eq_one_irreducible_all_ch12} collapses the peripheral set
    to $\{1\}$.
\end{proof}

\begin{remark}[Semigroup irreducible/primitive equivalence]\leanok
    The semigroup analog of the irreducible/primitive equivalence
    starts from one irreducible time slice, constructs a primitive
    smaller time step, and lifts the conclusion to the full semigroup
    (Theorems~\ref{thm:irreducible_semigroup_implies_primitive}
    and~\ref{thm:qds_irreducible_iff_primitive}).
\end{remark}

\begin{remark}[Reduction to a smaller time step for semigroups]\leanok
    Starting from one irreducible time $t_0$, we construct a
    positive time step $u=t_0/N$ with
    $T_{t_0}=(T_u)^N$, prove peripheral eigenvalue collapse at time~$u$,
    and deduce primitivity of $T_u$.
\end{remark}

\begin{theorem}[The stationary density is fixed for all times]
    \label{thm:fixed_density_fixed_for_all_times_ch12}
    \lean{fixed_density_fixed_for_all_times_of_irreducible_time}\leanok
    If $t_0>0$, $T_{t_0}$ is irreducible, and $\sigma$ is its unique fixed density
    matrix, then $T_u(\sigma)=\sigma$ for every $u\ge0$.
\end{theorem}

\begin{proof}
    \leanok
    By semigroup commutativity,
    $T_{t_0}(T_u(\sigma))=T_u(T_{t_0}(\sigma))=T_u(\sigma)$,
    and uniqueness of the fixed density at time $t_0$ identifies
    $T_u(\sigma)$ with $\sigma$.
\end{proof}

\begin{theorem}[Fixed points at the irreducible time are trace multiples]
    \label{thm:fixedPoint_eq_trace_smul_at_irreducible_time_ch12}
    \lean{fixedPoint_eq_trace_smul_at_irreducible_time}\leanok
    If $T_{t_0}$ is irreducible with faithful fixed density $\sigma$, then every fixed point
    $V$ of $T_{t_0}$ satisfies $V=\tr(V)\sigma$.
\end{theorem}

\begin{proof}
    \leanok
    This is the fixed-point trace-scalar theorem for the irreducible channel $T_{t_0}$.
\end{proof}

\begin{definition}[Residual slice index]
    \label{def:residual_slice_index_ch12}
    \lean{residualSliceIndex}\leanok
    For $u\ge0$ and $s>0$, the residual slice index is
    $m_n=\left\lfloor nu/s\right\rfloor$.
\end{definition}

\begin{definition}[Residual slice time]
    \label{def:residual_slice_time_ch12}
    \lean{residualSliceTime}\leanok
    For $u\ge0$ and $s>0$, the residual slice time is the remainder
    $r_n=s\,\mathrm{fract}\!\left(nu/s\right)$.
\end{definition}

\begin{theorem}[Residual slice decomposition]
    \label{thm:residual_slice_decomp_ch12}
    \lean{residualSliceTime_mem_Icc, residualSlice_decomp}\leanok
    \uses{def:residual_slice_index_ch12, def:residual_slice_time_ch12}
    For $u\ge0$ and $s>0$, one has $r_n\in[0,s]$ and $nu=m_ns+r_n$.
\end{theorem}

\begin{proof}
    \leanok
    This is the floor-plus-fraction decomposition of $nu/s$, multiplied by $s$.
\end{proof}

\begin{theorem}[A residual slice can vanish]
    \label{thm:exists_residual_time_eq_zero_ch12}
    \lean{exists_residual_time_eq_zero_of_fixedPoint}\leanok
    \uses{def:residual_slice_time_ch12, thm:residual_slice_decomp_ch12}
    Suppose $u\ge0$ and $s>0$, and that $T_s(\delta)=\delta$ and
    $T_{nu}(\delta)\to0$ as $n\to\infty$.
    Then there exists $a\in[0,s]$ such that $T_a(\delta)=0$.
\end{theorem}

\begin{proof}
    \leanok
    The residual slice times lie in the compact interval $[0,s]$, so a subsequence converges
    to some $a\in[0,s]$. The residual decomposition rewrites $T_{nu}(\delta)$ in terms of
    $T_{r_n}(\delta)$, and continuity of the semigroup identifies the limit with $T_a(\delta)$.
\end{proof}

\begin{theorem}[Reduced-time eigenvector has nonzero trace]
    \label{thm:exists_trace_ne_zero_eigenvector_of_fraction_slice_ch12}
    \lean{exists_trace_ne_zero_eigenvector_of_fraction_slice}\leanok
    Under the reduced-time hypotheses
    $T_{t_0}=(T_u)^{(\dim \MN{D})!}$, if
    $\mu$ is a peripheral eigenvalue of $T_u$, then $T_u$ admits a
    $\mu$-eigenvector with nonzero trace.
\end{theorem}

\begin{theorem}[Peripheral collapse at the reduced time step]
    \label{thm:peripheral_eq_one_of_fraction_slice_ch12}
    \lean{peripheral_eq_one_of_fraction_slice}\leanok
    Under the same hypotheses, every peripheral eigenvalue of $T_u$
    equals $1$.
\end{theorem}

\begin{theorem}[Primitivity at the reduced time step]
    \label{thm:primitive_of_fraction_slice_ch12}
    \lean{primitive_of_fraction_slice}\leanok
    Under the same hypotheses, $T_u$ is primitive.
\end{theorem}

\begin{theorem}[Irreducibility descends to a reduced time step]
    \label{thm:irreducible_of_fraction_slice_ch12}
    \lean{irreducible_of_fraction_slice}\leanok
    If $T_{t_0}=(T_u)^N$ and $T_{t_0}$ is irreducible, then $T_u$ is irreducible.
\end{theorem}

\begin{proof}
    \leanok
    Any invariant compression for $T_u$ remains invariant under powers, hence also for
    $T_{t_0}$; irreducibility of $T_{t_0}$ rules this out.
\end{proof}

\begin{theorem}[Existence of an irreducible reduced time step]
    \label{thm:exists_irreducible_fraction_slice_ch12}
    \lean{exists_irreducible_fraction_slice}\leanok
    If $t_0>0$, $T_{t_0}$ is irreducible, and $\sigma$ is fixed for all times,
    then there exists a positive time step $u$ such that
    $T_{t_0}=(T_u)^{(\dim \MN{D})!}$, the slice $T_u$ is a channel,
    $T_u$ is irreducible, and $T_u(\sigma)=\sigma$.
\end{theorem}

\begin{proof}
    \leanok
    \uses{thm:irreducible_of_fraction_slice_ch12}
    Take $u=t_0/N$ with $N=(\dim \MN{D})!$. The semigroup power identity gives the
    factorization of $T_{t_0}$, channelity comes from the semigroup hypotheses, and
    Theorem~\ref{thm:irreducible_of_fraction_slice_ch12} gives irreducibility of $T_u$.
\end{proof}

\begin{theorem}[Fixed points of the primitive reduced time step are trace-scalars]
    \label{thm:fixedPoint_eq_trace_smul_of_primitive_slice_ch12}
    \lean{fixedPoint_eq_trace_smul_of_primitive_slice}\leanok
    If $T_u$ is primitive and $\sigma$ is the common fixed density, then
    every fixed point of any positive-time slice $T_s$ has the form
    $\tau=\tr(\tau)\sigma$.
\end{theorem}

\begin{theorem}[Existence of a primitive reduced time step]
    \label{thm:exists_primitive_fraction_slice_ch12}
    \lean{exists_primitive_fraction_slice}\leanok
    If $T_{t_0}$ is irreducible for some $t_0>0$, then there exists a
    positive reduced time step $u$ such that $T_u$ is primitive.
\end{theorem}


\begin{theorem}[Irreducibility propagates to all positive times]
    \label{thm:irreducible_all_of_irreducible_time_ch12}
    \lean{irreducible_all_of_irreducible_time}\leanok
    \uses{def:is_quantum_dyn_semigroup_ch12,
        thm:fixed_density_fixed_for_all_times_ch12,
        thm:fixedPoint_eq_trace_smul_at_irreducible_time_ch12,
        thm:exists_primitive_fraction_slice_ch12,
        thm:fixedPoint_eq_trace_smul_of_primitive_slice_ch12}
    If $T_{t_0}$ is irreducible for some $t_0>0$, then every positive-time slice $T_s$ is
    irreducible.
\end{theorem}

\begin{proof}
    \leanok
    Let $\sigma$ be the unique faithful fixed density of $T_{t_0}$. By
    Theorem~\ref{thm:fixed_density_fixed_for_all_times_ch12}, it is fixed for all times, and
    Theorem~\ref{thm:fixedPoint_eq_trace_smul_at_irreducible_time_ch12} makes the fixed-point
    space of $T_{t_0}$ one-dimensional. Choose a primitive reduced time step by
    Theorem~\ref{thm:exists_primitive_fraction_slice_ch12}; then the trace-scalar description
    of fixed points from Theorem~\ref{thm:fixedPoint_eq_trace_smul_of_primitive_slice_ch12}
    gives the uniqueness criterion for irreducibility at every positive time.
\end{proof}

\begin{theorem}[An irreducible semigroup is primitive]
    \label{thm:irreducible_semigroup_implies_primitive}
    \lean{irreducible_semigroup_implies_primitive}\leanok
    \uses{def:is_quantum_dyn_semigroup_ch12,
        thm:irreducible_all_of_irreducible_time_ch12,
        thm:primitive_of_irreducible_all_ch12}
    Let $T_t=e^{tL}$ be a quantum dynamical semigroup.
    If $T_{t_0}$ is irreducible for some $t_0>0$, then
    $T_t$ is primitive (and irreducible) for every $t>0$.
    This is \cite[Proposition~7.5]{Wolf2012Quantum}.
\end{theorem}

\begin{proof}\leanok
    First apply Theorem~\ref{thm:irreducible_all_of_irreducible_time_ch12} to propagate
    irreducibility from the single time $t_0$ to every positive time. Then apply
    Theorem~\ref{thm:primitive_of_irreducible_all_ch12} to conclude primitivity of every
    positive-time slice.
\end{proof}

\begin{theorem}[Irreducibility iff primitivity for a quantum dynamical semigroup]
    \label{thm:qds_irreducible_iff_primitive}
    \lean{qds_irreducible_iff_primitive}\leanok
    \uses{thm:irreducible_semigroup_implies_primitive}
    For a quantum dynamical semigroup $T_t=e^{tL}$,
    \begin{align}
        (\exists\, t_0>0,\ T_{t_0}\ \text{irreducible})
        &\iff
        (\forall\, t>0,\ T_t\ \text{primitive and irreducible}).
        \label{eq:semigroup_qds_equiv}
    \end{align}
    This is \cite[Proposition~7.5]{Wolf2012Quantum}.
\end{theorem}

\begin{proof}\leanok
    \uses{thm:irreducible_semigroup_implies_primitive}
    Forward: apply Theorem~\ref{thm:irreducible_semigroup_implies_primitive}.
    Reverse: take $t_0=1$ (or any positive time); the right-hand side
    already includes irreducibility of $T_t$ for every $t>0$.
\end{proof}

%% ---------------------------------------------------------
